\documentclass[11pt]{article}

\usepackage[T1]{fontenc}
\usepackage{lmodern}
\usepackage{amsmath,amssymb,amsthm}
\usepackage[margin=1in]{geometry}
\usepackage{booktabs}
\usepackage{microtype}
\usepackage[hidelinks]{hyperref}

\newtheorem{theorem}{Theorem}
\newtheorem{proposition}[theorem]{Proposition}
\newtheorem{lemma}[theorem]{Lemma}
\newtheorem{corollary}[theorem]{Corollary}
\theoremstyle{remark}
\newtheorem{remark}[theorem]{Remark}

\newcommand{\rad}{\operatorname{rad}}
\newcommand{\ord}{\operatorname{ord}}
\newcommand{\QQ}{\mathbb{Q}}
\newcommand{\OK}{\mathcal{O}_K}

\title{An iterated quadratic \(abc\)-orbit and a Lucas--Wieferich obstruction}
\author{Magnus Gille\\Independent researcher}
\date{July 2026}

\begin{document}
\maketitle

\begin{abstract}
The elementary transfer
\[
 (a,b,a+b)\longmapsto\bigl(4ab,(a-b)^2,(a+b)^2\bigr)
\]
is iterated from \((1,8,9)\).  Writing the resulting triples as
\((a_n,b_n,c_n)\), we prove the exact radical formula
\[
 \frac{\rad(a_nb_nc_n)}{c_n}
 =\frac{2}{3}\,
   \frac{\prod_{j<n}|T_{2^j}(7/9)|}
        {\prod_{j<n}|d_j|/\rad(d_j)},\qquad d_j=a_j-b_j.
\]
In particular every member of the orbit is an \(abc\)-hit.  An explicit
application of the Baker--Wüstholz theorem shows that the logarithm of the
Chebyshev product is \(O(n^2)\), whereas \(\log c_n=2^n\log 9\).  Thus the
asymptotic \(abc\)-quality of this orbit is governed, up to \(O(n^2)\), only
by repeated prime factors of the pairwise coprime integers \(d_j\).
These integers form the dyadic subsequence
\[
 |d_j|=\frac12\bigl|V_{2^{j+1}}(2,9)\bigr|
\]
of a Lucas \(V\)-sequence.  A prime divisor has prescribed dyadic order in
\(\QQ(\sqrt{-2})\), and a square divisor is exactly a Wieferich-type failure
of order growth.  A reproducible search found no such square divisor for
primes \(p\leq 10^7\) and indices \(0\leq j\leq50\).  The computation is
not used in the proofs.
\end{abstract}

\section{Introduction}

For a positive integer \(m\), let
\(\rad(m)=\prod_{p\mid m}p\).  A coprime triple of positive integers
\((a,b,c)\) with \(a+b=c\) is called an \(abc\)-hit when
\(\rad(abc)<c\), and its quality is
\[
 q(a,b,c)=\frac{\log c}{\log\rad(abc)}.
\]
The \(abc\) conjecture predicts that, for every \(\varepsilon>0\), only
finitely many such triples satisfy \(c>\rad(abc)^{1+\varepsilon}\);
see, for example, Masser's formulation~\cite{Masser1985}.

The identity
\[
 4ab+(a-b)^2=(a+b)^2                                      \tag{1.1}
\]
transfers one additive triple to another.  Polynomial transfers of this
kind have long been used in searches for \(abc\)-triples; see van der
Horst~\cite[\S2.3]{vanderHorst2010}.  Neither (1.1) nor the linear
recurrence appearing below is new.  In particular, the sequence
\(3^mT_m(1/3)=V_m(2,9)/2\) is OEIS A025172~\cite{OEIS-A025172}.

This note analyzes what happens when (1.1) is iterated from the small hit
\((1,8,9)\).  Its contribution is an exact separation of the radical loss
into an archimedean Chebyshev factor and a nonarchimedean repeated-prime
factor.  A quantitative linear-form estimate makes the first factor
sub-power-sized.  The remaining question is then a radical problem for a
fixed Lucas sequence at dyadic indices.  We also give an elementary order
criterion which identifies square factors as number-field
Lucas--Wieferich phenomena.

The conclusions are deliberately limited.  We prove an infinite family of
\(abc\)-hits, but no fixed quality gap.  In particular, the results neither
prove nor disprove the \(abc\) conjecture.

\section{The quadratic orbit}

Set
\[
 (a_0,b_0,c_0)=(1,8,9)
\]
and, for \(n\geq0\), define
\[
 (a_{n+1},b_{n+1},c_{n+1})
   =\bigl(4a_nb_n,(a_n-b_n)^2,c_n^2\bigr).                 \tag{2.1}
\]
Identity (1.1) gives \(a_n+b_n=c_n\) at every index.  Put
\[
 d_n=a_n-b_n.
\]
Then
\[
 c_n=9^{2^n},\qquad
 d_{n+1}=c_n^2-2d_n^2.                                    \tag{2.2}
\]

\begin{lemma}\label{lem:coprime}
Each \(d_n\) is odd and coprime to \(6\), and the integers \(d_n\) are
pairwise coprime.
\end{lemma}

\begin{proof}
The assertions about \(2\) and \(3\) follow from \(d_0=-7\) and (2.2):
all \(c_n\) are odd and divisible by \(3\), while
\(d_{n+1}\equiv d_n^2\equiv1\pmod3\).

Suppose that a prime \(p\) divides \(d_m\).  Since \(p\nmid c_m\), (2.2)
gives
\[
 d_{m+1}\equiv c_{m+1}\pmod p,\qquad
 d_{m+2}\equiv-c_{m+2}\pmod p.
\]
If \(d_k\equiv-c_k\pmod p\), then (2.2) gives
\(d_{k+1}\equiv-c_{k+1}\pmod p\).  Hence \(p\) divides no \(d_k\) with
\(k>m\).  Interchanging two indices proves pairwise coprimality.
\end{proof}

\begin{lemma}\label{lem:support}
Every triple in (2.1) is primitive and
\[
 \rad(a_nb_nc_n)=6\prod_{j=0}^{n-1}\rad(|d_j|),            \tag{2.3}
\]
where an empty product is \(1\).
\end{lemma}

\begin{proof}
For \(n=0\), the prime support of \(a_0b_0\) is \(\{2\}\).  For \(n\geq1\),
\(b_n=d_{n-1}^2\), while the recurrence for \(a_n\) in (2.1) shows
inductively that the prime support of \(a_n\) is
\[
 \{2\}\cup\bigcup_{j=0}^{n-2}\{p:p\mid d_j\}.
\]
Lemma~\ref{lem:coprime} makes these supports disjoint and excludes \(3\).
The only prime in \(c_n\) is \(3\), which proves both claims.
\end{proof}

Normalize (2.2) by setting \(x_n=-d_n/c_n\).  Since
\(T_2(x)=2x^2-1\),
\[
 x_{n+1}=T_2(x_n),\qquad x_0=\frac79,
\]
and therefore
\[
 -\frac{d_n}{c_n}=T_{2^n}\left(\frac79\right).             \tag{2.4}
\]

\section{The exact radical identity}

Define
\[
 t_j=\frac{|d_j|}{c_j}
     =\left|T_{2^j}\left(\frac79\right)\right|,
 \qquad
 Q_n=\prod_{j=0}^{n-1}\frac{|d_j|}{\rad(|d_j|)}.           \tag{3.1}
\]
Thus \(Q_n\) records exactly the prime multiplicities beyond the first
within \(d_0,\ldots,d_{n-1}\).

\begin{theorem}[Exact separation]\label{thm:separation}
For every \(n\geq0\),
\[
 \boxed{\quad
 \frac{\rad(a_nb_nc_n)}{c_n}
 =\frac23\,\frac{\prod_{j=0}^{n-1}t_j}{Q_n}.
 \quad}                                                    \tag{3.2}
\]
In particular, \((a_n,b_n,c_n)\) is an \(abc\)-hit for every \(n\).
\end{theorem}

\begin{proof}
As \(\prod_{j<n}c_j=9^{2^n-1}=c_n/9\), (3.1) gives
\[
 \prod_{j<n}\rad(|d_j|)
 =\frac{\prod_{j<n}|d_j|}{Q_n}
 =\frac{c_n}{9}\frac{\prod_{j<n}t_j}{Q_n}.
\]
Substitution into (2.3) proves (3.2).  For \(n=0\), its right side is
\(2/3\).  For \(n>0\), each \(t_j\) lies strictly between \(0\) and \(1\);
this also follows from the non-root-of-unity argument in the proof of
Theorem~\ref{thm:baker}.  Since \(Q_n\geq1\), the ratio in (3.2) is
strictly less than \(1\).
\end{proof}

\section{The archimedean factor}

The estimate in this section is stated with an explicit, if very large,
constant so that the dependence on the index is unambiguous.

\begin{theorem}[Archimedean factor]\label{thm:baker}
Let
\[
 C_{\mathrm{BW}}
 =18\cdot3!\cdot2^3\cdot64^4\log 8,
 \qquad
 \kappa=C_{\mathrm{BW}}(\log3)\frac{\pi}{2}.
\]
Then, for \(j\geq0\),
\[
 t_j>\frac1\pi\,2^{-\kappa(j+1)}.                          \tag{4.1}
\]
Consequently,
\[
 0<-\sum_{j=0}^{n-1}\log t_j
 <n\log\pi+\frac{\kappa\log2}{2}n(n+1).                   \tag{4.2}
\]
\end{theorem}

\begin{proof}
Put
\[
 z=\frac{7+4\sqrt{-2}}9=e^{i\theta},
 \qquad \theta=\arccos(7/9)\in(0,\pi/2).
\]
Then \(z+z^{-1}=14/9\).  The number \(z\) is not a root of unity:
otherwise \(z+z^{-1}\) would be both a rational number and an algebraic
integer, hence an integer.

For an integer \(N\geq1\), choose an odd integer \(k\) nearest to
\(2N\theta/\pi\), and use the principal logarithms
\(\log z=i\theta\) and \(\log(-1)=i\pi\).  The linear form
\[
 \Lambda_N=2N\log z-k\log(-1)                              \tag{4.3}
\]
is nonzero, satisfies \(|\Lambda_N|\leq\pi\), and has
\(|k|\leq N+1\leq2N\).

We apply the explicit Baker--Wüstholz estimate
\cite[Theorem~1]{BakerWustholz1993}.  In its standard notation, the number
field \(\QQ(z,-1)=\QQ(\sqrt{-2})\) has degree \(D=2\), there are two
logarithms, and the coefficient bound may be taken as \(B=2N\).  The
minimal polynomial \(9X^2-14X+9\) has both roots on the unit circle, so
\[
 h(z)=\log3,\qquad
 \max\{h(z),|\log z|/D,1/D\}=\log3.
\]
For \(-1\), the corresponding modified height is \(\pi/2\).  The theorem
therefore gives
\[
 \log|\Lambda_N|>
 -C_{\mathrm{BW}}(\log3)\frac{\pi}{2}\log(2N)
 =-\kappa\log(2N).                                        \tag{4.4}
\]

Because \(k\) is odd,
\[
 |z^{2N}+1|=|1-e^{\Lambda_N}|.
\]
For a purely imaginary \(\Lambda\) with \(|\Lambda|\leq\pi\),
\(|1-e^\Lambda|\geq2|\Lambda|/\pi\).  Hence (4.4) yields
\[
 \frac12|z^{2N}+1|>\frac1\pi(2N)^{-\kappa}.                \tag{4.5}
\]
Taking \(N=2^j\) and using
\[
 t_j=\frac12|z^{2N}+1|
\]
proves (4.1).  Summing its logarithm proves the upper bound in (4.2).
The lower bound follows from \(0<t_j<1\); equality \(t_j=1\) would make
\(z\) a root of unity.
\end{proof}

\begin{corollary}[Quality reduction]\label{cor:quality}
Let
\[
 R_n=\rad(a_nb_nc_n),\quad L_n=\log c_n=2^n\log9,\quad
 q_n=\frac{\log c_n}{\log R_n}.
\]
Then
\[
 \log\frac{c_n}{R_n}=\log Q_n+E_n,                         \tag{4.6}
\]
where
\[
 \log\frac32\leq E_n
 <\log\frac32+n\log\pi+
   \frac{\kappa\log2}{2}n(n+1).                            \tag{4.7}
\]
In particular,
\[
 \log\frac{c_n}{R_n}=\log Q_n+O(n^2),                     \tag{4.8}
\]
and
\[
 q_n\longrightarrow1
 \quad\Longleftrightarrow\quad
 \log Q_n=o(\log c_n).                                     \tag{4.9}
\]
More generally,
\[
 \limsup_{n\to\infty}(q_n-1)>0
 \quad\Longleftrightarrow\quad
 \limsup_{n\to\infty}\frac{\log Q_n}{\log c_n}>0.          \tag{4.10}
\]
\end{corollary}

\begin{proof}
Taking logarithms in (3.2) gives (4.6) with
\[
 E_n=\log(3/2)-\sum_{j<n}\log t_j.
\]
Theorem~\ref{thm:baker} gives (4.7)--(4.8).  Since \(n^2=o(L_n)\),
(4.9) follows from
\[
 q_n-1=\frac{\log(c_n/R_n)}{\log R_n}.
\]
The same identity and (4.8) show that either side of (4.10) is positive
exactly when the other is.
\end{proof}

\begin{remark}
For a fixed \(\delta>0\), the exact inequality
\[
 q_n\geq1+\delta
 \quad\Longleftrightarrow\quad
 \log Q_n+E_n\geq\frac{\delta}{1+\delta}\log c_n
\]
is sometimes a more useful form of the reduction.
\end{remark}

\section{Lucas sequence and order lifting}

Let \(V_m(P,Q)\) denote the Lucas \(V\)-sequence with
\[
 V_0=2,\quad V_1=P,\quad V_m=PV_{m-1}-QV_{m-2}.
\]
For \(P=2,Q=9\), put
\[
 A_m=\frac12V_m(2,9)=3^mT_m(1/3).
\]
Thus \(A_0=A_1=1\) and \(A_m=2A_{m-1}-9A_{m-2}\).  This is the sequence
OEIS A025172~\cite{OEIS-A025172}.

\begin{proposition}\label{prop:lucas}
For every \(j\geq0\),
\[
 |d_j|=|A_{2^{j+1}}|
      =\frac12\bigl|V_{2^{j+1}}(2,9)\bigr|.                \tag{5.1}
\]
\end{proposition}

\begin{proof}
The composition rule \(T_r(T_s(x))=T_{rs}(x)\) and
\(T_2(1/3)=-7/9\) give
\[
 A_{2^{j+1}}
 =9^{2^j}T_{2^j}(-7/9).
\]
Comparing with (2.4), and using \(T_m(-x)=(-1)^mT_m(x)\), proves (5.1).
\end{proof}

The next proposition locates every prime factor in a thin congruence
class and makes the Wieferich description precise.  Put
\[
 K=\QQ(\sqrt{-2}),\qquad
 \alpha=1+2\sqrt{-2},\quad\beta=1-2\sqrt{-2},\quad
 u=\alpha/\beta.
\]
Then \(V_m(2,9)=\alpha^m+\beta^m\).

\begin{proposition}[Dyadic order]\label{prop:order}
Let \(p\) be a prime divisor of \(d_j\), put \(m=2^{j+1}\), and let
\(\mathfrak p\) be a prime of \(\OK\) above \(p\).  Then \(p\nmid6\) and
\[
 \ord_{(\OK/\mathfrak p)^\times}(u)=2m=2^{j+2}.            \tag{5.2}
\]
Consequently,
\[
 \begin{cases}
 p\equiv 1\pmod{2^{j+2}},&p\text{ splits in }K,\\
 p\equiv-1\pmod{2^{j+2}},&p\text{ is inert in }K.
 \end{cases}                                               \tag{5.3}
\]
Moreover,
\[
 p^2\mid d_j
 \quad\Longleftrightarrow\quad
 u^m\equiv-1\pmod{\mathfrak p^2}
 \text{ for every }\mathfrak p\mid p,                     \tag{5.4}
\]
and in that case the order in (5.2) does not acquire the usual factor
\(p\) on lifting from \(\mathfrak p\) to \(\mathfrak p^2\).
\end{proposition}

\begin{proof}
Lemma~\ref{lem:coprime} gives \(p\nmid6\).  Since \(\beta\) has norm \(9\),
it is invertible modulo \(\mathfrak p\).  From (5.1),
\(\alpha^m+\beta^m\equiv0\pmod{\mathfrak p}\), and hence
\(u^m\equiv-1\pmod{\mathfrak p}\).  As \(m\) is a power of \(2\), the
order of \(u\) is exactly \(2m\).

If \(p\) splits, the residue field has \(p\) elements, so \(2m\mid p-1\).
If \(p\) is inert, conjugation is the Frobenius on the residue field and
\(u^{p+1}=u\bar u=1\), so \(2m\mid p+1\).  This proves (5.3).

Because \(p\) is unramified and \(2\beta^m\) is a unit at every
\(\mathfrak p\mid p\), (5.1) gives (5.4).  The reduction of \(u\) has
order prime to \(p\); its order modulo \(\mathfrak p^2\) is therefore
either \(2m\) or \(2mp\).  In the first case \(u^m\) is the unique lift of
the element \(-1\) of order \(2\), so it is \(-1\) modulo
\(\mathfrak p^2\).  The converse is immediate.
\end{proof}

Thus a prime with \(p^2\mid d_j\) is naturally a Lucas--Wieferich prime
for the algebraic number \(u\), which is a unit away from \(3\).
Classical primitive-divisor theory for
Lucas and Lehmer sequences addresses the occurrence of new prime divisors
rather than their aggregate multiplicities; see, for example,
Bilu--Hanrot--Voutier~\cite{BiluHanrotVoutier2001}.  In this orbit,
Lemma~\ref{lem:coprime} already makes every prime divisor new relative to
all other dyadic terms.  What (4.9)--(4.10) require instead is control of
the total repeated-prime quotient \(Q_n\).  Related squarefree primitive
divisor results in arithmetic dynamics are conditional on Vojta's
conjecture over number fields~\cite{GhiocaNguyenTucker2018}, and in any
case the existence of one squarefree prime per term would be much weaker
than the proportional radical estimate needed here.

\section{Computation}

The accompanying program \texttt{check\_square\_lifts.py} works entirely
modulo \(p^2\).  For each prime \(p\), it starts with
\((c,d)=(9,-7)\) and iterates
\[
 (c,d)\longmapsto(c^2,c^2-2d^2)\pmod{p^2}.
\]
It reports \(p,j\) precisely when \(d\equiv0\pmod{p^2}\).  If
\(d\equiv0\pmod p\) but not modulo \(p^2\), the search for that prime may
stop, by Lemma~\ref{lem:coprime}.

\begin{proposition}[Finite verification]\label{prop:computation}
For every prime \(p\leq10^7\), \(p\neq2,3\), and every
\(0\leq j\leq50\), the modular recurrence gives
\[
 p^2\nmid d_j.
\]
The search tests \(664{,}577\) primes and reports zero square lifts.
\end{proposition}

This is a finite computer verification, not an input to any preceding
theorem.  The script uses a deterministic sieve and integer modular
arithmetic only.  Its command line is
\[
\texttt{python3 check\_square\_lifts.py --prime-limit 10000000 --max-j 50}.
\]
For orientation, the first terms factor as follows.
\[
\begin{array}{c@{\qquad}r@{\qquad}l}
\toprule
j&d_j&|d_j|\text{ factored}\\
\midrule
0&-7&7\\
1&-17&17\\
2&5983&31\cdot193\\
3&-28545857&2753\cdot10369\\
4&223288285122943&127\cdot19841\cdot88613249\\
\bottomrule
\end{array}
\]

\section{Concluding question}

The orbit has an exact dichotomy.  Its quality tends to \(1\) if and only
if the accumulated repeated-prime quotient of the dyadic Lucas terms is
sub-power-sized:
\[
 \prod_{j<n}\frac{|V_{2^{j+1}}(2,9)|/2}
                    {\rad(|V_{2^{j+1}}(2,9)|/2)}
   =c_n^{o(1)}.
\]
A fixed quality gap occurs along a subsequence if and only if the same
product has positive-power growth.  The Baker--Wüstholz estimate rules out
the archimedean Chebyshev values as a source of such a gap; the remaining
obstruction is purely the aggregate of Lucas--Wieferich lifts.  We do not
know how to prove either alternative.

\section*{AI-use statement}

Claude Fable 5 and OpenAI Codex were used extensively to generate candidate
arguments, conduct literature searches, write and run the computation,
draft the manuscript, and perform reciprocal adversarial review.  The
named author directed the work and is responsible for the final content.

\begin{thebibliography}{99}

\bibitem{BakerWustholz1993}
A.~Baker and G.~Wüstholz,
\newblock Logarithmic forms and group varieties,
\newblock \emph{J. Reine Angew. Math.} \textbf{442} (1993), 19--62.
\newblock \url{https://doi.org/10.1515/crll.1993.442.19}.

\bibitem{BiluHanrotVoutier2001}
Y.~Bilu, G.~Hanrot, and P.~M.~Voutier,
\newblock Existence of primitive divisors of Lucas and Lehmer numbers,
\newblock \emph{J. Reine Angew. Math.} \textbf{539} (2001), 75--122.
\newblock \url{https://doi.org/10.1515/crll.2001.080}.

\bibitem{GhiocaNguyenTucker2018}
D.~Ghioca, K.~D.~Nguyen, and T.~J.~Tucker,
\newblock Squarefree doubly primitive divisors in dynamical sequences,
\newblock \emph{Math. Proc. Cambridge Philos. Soc.} \textbf{164} (2018),
551--572.
\newblock \url{https://doi.org/10.1017/S0305004117000500}.

\bibitem{Masser1985}
D.~W.~Masser,
\newblock Open problems,
\newblock in \emph{Proceedings of the Symposium on Analytic Number Theory},
W.~W.~L.~Chen, ed., Imperial College, London, 1985, pp.~591--607.

\bibitem{OEIS-A025172}
N.~J.~A.~Sloane and The OEIS Foundation,
\newblock The On-Line Encyclopedia of Integer Sequences, A025172,
\newblock \url{https://oeis.org/A025172}, accessed July 26, 2026.

\bibitem{vanderHorst2010}
J.~P.~van der Horst,
\newblock \emph{Finding \(ABC\)-triples using elliptic curves},
\newblock M.Sc. thesis, Universiteit Leiden, 2010.

\end{thebibliography}

\end{document}
